% secoes/kl_cilindro.tex -- KL de Galerkin do NeoPZ no meio cilindro de Mika et al. (2021), Secao 6.2, Exemplo 1-1,
% e a conferencia do notebook randomfied-v2.nb. Figuras e tabelas: Projects/KLCilindroMika/figuras/cilindro.py;
% resultados e procedimento: Projects/KLCilindroMika/README.md.
\section{Karhunen-Loève no meio cilindro de Mika et al. (2021)}\label{sec:kl:cil}

\subsection{O problema}\label{sec:kl:cil:problema}

O Exemplo 1-1 da Seção~6.2 de Mika, Hughes, Schillinger, Wriggers e Hiemstra (2021, CMAME 379:113730) é um campo
aleatório com variância $\sigma^2 = 1$ e núcleo exponencial isotrópico
\begin{equation}
  C(\mathbf{x}, \mathbf{x}') = \exp\left(-\frac{\lVert \mathbf{x} - \mathbf{x}' \rVert}{5}\right)
  \label{eq:kl:cil:nucleo}
\end{equation}
(distância euclidiana, comprimento de correlação de 5~m) no meio cilindro vazado
$D = \{8 \le \rho \le 10,\ 0 \le \theta \le \pi,\ 0 \le z \le 15\}$ (m), de volume $|D| = 270\pi = 848{,}2300$~m$^3$
(Figura~\ref{fig:kl:cil:geometria}). Com $\sigma^2 = 1$, a soma de todos os autovalores exatos é
$\int_D C(\mathbf{x}, \mathbf{x})\,d\mathbf{x} = 270\pi$. O exemplo é útil como teste da KL de Galerkin do NeoPZ em
3D com geometria curva, porque o artigo publica os 20 maiores autovalores (Tabela~4) com o custo de cada cálculo. O
programa é o \texttt{KLCilindroMika} (\texttt{Projects/KLCilindroMika/main.cpp}); os números desta seção estão no
README desse diretório e em \texttt{dados/kl\_cil\_*.csv}.

\begin{figure}[htbp]
  \centering
  \includegraphics[width=0.6\linewidth]{kl_cil_geometria.pdf}
  \caption{O meio cilindro vazado de Mika et al.\ (2021) com as arestas da malha $12 \times 2 \times 6$
    ($n_\theta \times n_\rho \times n_z$, 144 hexaedros) na superfície, desenhadas pelo mapa cilíndrico exato.}
  \label{fig:kl:cil:geometria}
\end{figure}

\subsection{As formulações de Mika et al.}\label{sec:kl:cil:mika}

O artigo propõe um método de Galerkin isogeométrico sem montagem de matrizes (\emph{matrix-free}) para a KL, com três
ingredientes: (1) uma quadratura baseada em interpolação (IBQ): o núcleo é interpolado num espaço de \emph{splines}
produto tensorial (o ``espaço de interpolação''), o que reduz a integral dupla a contrações tensoriais com matrizes
univariadas e minimiza o número de pontos de quadratura; (2) uma reformulação barata do problema generalizado num
problema padrão; (3) um produto matriz-vetor sem matriz, paralelo, dentro do Lanczos. A geometria é representada
exatamente por dois \emph{patches} NURBS (graus $\{2, 1, 1\}$, vetores de nós
$\Xi_1 = (0, 0, 0, 0{,}5, 0{,}5, 1, 1, 1)$, $\Xi_2 = \Xi_3 = (0, 0, 1, 1)$), com continuidade $C^0$ entre os dois.
A Tabela~4 do artigo traz três casos da IBQ (espaço de solução de 256 elementos e 1050 graus de liberdade; espaços de
interpolação de 1980, 8990 e 16\,770 graus de liberdade) e duas soluções de referência pelo Galerkin isogeométrico
padrão de Rahman (2018, CMAME 338:533--561), com o núcleo exato, espaços NURBS de grau $(2, 2, 2)$ e montagem elemento a elemento com
quadratura de Gauss completa:
\begin{itemize}
  \item \textbf{Galerkin 1}: 256 elementos, 1050 graus de liberdade, $h = 2{,}857$; formação e montagem 4,86~min,
    solução 0,09~s;
  \item \textbf{Galerkin 2}: 3800 elementos, 6642 graus de liberdade, $h = 1{,}073$; formação e montagem 17,3~h,
    solução (Lanczos, 63 iterações) 6,21~s, 1,566~GB; um núcleo de processador.
\end{itemize}
Galerkin 2 é a coluna mais precisa da tabela. Aqui as duas são comparadas com uma referência de Nyström convergida
(Seção~\ref{sec:kl:cil:nystrom}).

\subsection{A KL de Galerkin 3D do NeoPZ}\label{sec:kl:cil:galerkin}

A formulação é a da Seção~\ref{kl:form:galerkin}: $\mathbf{C}\mathbf{a} = \lambda \mathbf{B}\mathbf{a}$ com
$\mathbf{B}$ e $\mathbf{C}$ densas. As particularidades do caso 3D são:
\begin{itemize}
  \item \textbf{Espaço e geometria}: hexaedros de ordem 2 (Q2, 27 funções por elemento, $H^1$ contínuo) numa malha
    estruturada $n_\theta \times n_\rho \times n_z$; a geometria é exata (\texttt{TPZGeoBlend} com \texttt{TPZArc3D}
    nas arestas circunferenciais): a soma dos pesos vezes $|\det J|$ dá $270\pi$ com erro relativo
    $1{,}5 \times 10^{-15}$, e os pontos mapeados diferem do mapa cilíndrico exato em até $5{,}3 \times 10^{-15}$~m
    (malha $4 \times 1 \times 2$). As malhas são $8 \times 1 \times 4$, $12 \times 2 \times 6$, $16 \times 2 \times 8$ e
    $24 \times 3 \times 12$ (32 a 864 elementos, 459 a 8575 equações).
  \item \textbf{Quadratura}: $\mathbf{B}$ com a regra do elemento; em $\mathbf{C}$, cada par de elementos
    $(e_1, e_2)$ é integrado com a regra de Gauss produto de cada elemento, com $3^3$ pontos (a regra do elemento,
    $P + 1$ pontos por direção) ou com $5^3$ pontos em todos os pares (\texttt{NG=5}): são $27^2$ ou $125^2$
    avaliações do núcleo por par. O núcleo~\eqref{eq:kl:cil:nucleo} não é derivável na diagonal
    $\mathbf{x} = \mathbf{x}'$, que fica nos pares de um elemento com ele mesmo (e toca os vizinhos nas faces comuns);
    a regra de Gauss não integra esses pares com a precisão da sua ordem.
  \item \textbf{Montagem e solução}: só os pares $e_1 \le e_2$ (simetria), com 2 threads; os autovalores saem do
    Eigen (Cholesky de $\mathbf{B}$ e \texttt{SelfAdjointEigenSolver}, todos os $N$ autovalores, 1 thread). O custo é
    $O(N^2)$ em memória e $O(N^3)$ em tempo: com 8575 equações, 337~s no problema de autovalores e 1,7~GB; a montagem de
    $\mathbf{C}$ levou 6,5~s com $3^3$ pontos e 78~s com $5^3$.
\end{itemize}
A conferência \texttt{VALIDAR=1} (malha $4 \times 1 \times 2$) compara o \texttt{CalcStiffC} corrigido com o
algoritmo anterior em pares de elementos iguais ($1{,}4 \times 10^{-15}$), a simetria
$\mathbf{C}(e_1, e_2) = \mathbf{C}(e_2, e_1)^T$ ($3{,}4 \times 10^{-16}$) e a matriz montada pelos dados dos pontos
com \texttt{CalcStiffC} par a par ($1{,}2 \times 10^{-16}$).

\paragraph{Defeitos corrigidos na biblioteca.}
Em \texttt{TPZKarhunenLoeveMat}: \texttt{Dimension()} devolvia 2 para qualquer malha; o núcleo gaussiano (tipo 2) não
tinha \texttt{break} e devolvia o exponencial separável; um tipo de núcleo não definido devolvia 0; construtores
deixavam $L_x$, $L_y$, $L_z$, o tipo e a dimensão sem valor; \texttt{NSolutionVariables}, \texttt{SetData} e
\texttt{SetLinearContext} chamavam a si mesmos; \texttt{Solution} sobrescrevia \texttt{data.sol[0]}. Em
\texttt{TPZInterpolationSpace::CalcStiffC}: a regra de integração do primeiro elemento era usada também no segundo,
a matriz do par tinha $n_\text{shape}(e_1)$ colunas em vez de $n_\text{shape}(e_2)$, e os dados do segundo elemento
eram recalculados para cada par de pontos ($n_{p1} n_{p2}$ chamadas, agora $n_{p1} + n_{p2}$). Foram acrescentados
\texttt{KLPointData}, \texttt{CalcStiffC} com uma regra por elemento e \texttt{ContributeC} a partir dos dados dos
pontos.

\subsection{Referência de Nyström}\label{sec:kl:cil:nystrom}

A referência é o método de Nyström da Seção~\ref{kl:form:nystrom}, em numpy/scipy, sem o NeoPZ: Gauss-Legendre
composto em $(\rho, \theta, z)$ com pesos $\rho\,d\rho\,d\theta\,dz$, na forma simétrica
$\mathbf{W}^{1/2}\mathbf{K}\mathbf{W}^{1/2}$. As simetrias $\theta \to \pi - \theta$ e $z \to 15 - z$ separam o
problema em 4 classes, cada uma no quarto $0 \le \theta \le \pi/2$, $0 \le z \le 7{,}5$ com o núcleo somado nas
imagens; os maiores autovalores de cada classe saem do Lanczos (\texttt{eigsh}). O quarto tem
$m \times 7m \times 4m$ painéis em $\rho \times \theta \times z$ com $q$ pontos por painel e direção: $m = 1$ a 4 com
$q = 2$ (até 14\,336 pontos por classe), $m = 1$ a 3 com $q = 3$ (até 20\,412) e $m = 1$ e 2 com $q = 4$. Os
autovalores convergem com $h^4$ ($q = 2$, $m = 2, 3, 4$: as diferenças sucessivas de $\lambda_1$ caem 5,80 vezes;
$h^4$ dá 5,94). A referência é a média de quatro extrapolações de Richardson ($h^4 + h^5$ e $h^4 + h^6$, com $q = 2$
em $m = 2, 3, 4$ e com $q = 3$ em $m = 1, 2, 3$); as quatro diferem em no máximo $4{,}2 \times 10^{-7}$ do valor
($\lambda_1$: $5{,}9 \times 10^{-9}$), e a soma dos traços das 4 classes é $270\pi$. Resultado:
$\lambda_1 = 162{,}796422$, \ldots, $\lambda_{20} = 5{,}761164$ (Tabela~\ref{tab:kl:cil:autovalores}).

\paragraph{Por que não há forma fechada.}
No talude, o núcleo é separável, $\exp(-|x - x'|/L_x - |y - y'|/L_y)$: é o produto de dois núcleos exponenciais 1D,
cada um com KL em forma fechada no intervalo (a equação integral vira uma EDO de segunda ordem com condições de Robin),
e no retângulo os autopares 2D são os produtos dos 1D (Seção~\ref{kl:form:analitica}). O núcleo
\eqref{eq:kl:cil:nucleo} usa a distância euclidiana e não é produto de núcleos 1D em nenhum sistema de coordenadas: em
coordenadas cilíndricas,
$\lVert \mathbf{x} - \mathbf{x}' \rVert^2 = \rho^2 + \rho'^2 - 2\rho\rho'\cos(\theta - \theta') + (z - z')^2$, que
não separa as variáveis. Em $\mathbb{R}^3$, $e^{-r/b}$ é, a menos de constante, a função de Green de
$(1 - b^2\Delta)^2$, mas restrito a um domínio limitado as condições de contorno equivalentes são não locais; e o meio
anel ($0 \le \theta \le \pi$) não tem a invariância por rotação que daria modos de Fourier em $\theta$. Por isso a
referência é numérica.

\subsection{Resultados}\label{sec:kl:cil:resultados}

A Tabela~\ref{tab:kl:cil:convergencia} e a Figura~\ref{fig:kl:cil:convergencia} mostram o erro relativo dos
autovalores do NeoPZ contra a referência de Nyström; a Tabela~\ref{tab:kl:cil:autovalores} e a
Figura~\ref{fig:kl:cil:espectro}, os 20 autovalores.
\begin{itemize}
  \item Nas oito contas, os 20 autovalores do NeoPZ ficam acima da referência. Como, com integração exata, os
    autovalores de Galerkin seriam limites inferiores dos exatos (princípio do min-max), o erro medido vem da
    quadratura de $\mathbf{C}$. Com $5^3$ pontos, o erro cai, na mesma malha, 6,0 a 6,3 vezes em $\lambda_1$ e 6,7 a
    8,7 vezes no máximo dos 20 modos; a malha $16 \times 2 \times 8$ com $5^3$ pontos (2805 equações,
    $1{,}6 \times 10^{-4}$) já fica abaixo da $24 \times 3 \times 12$ com $3^3$ ($2{,}2 \times 10^{-4}$).
  \item O erro cai com $h^4$: de $16 \times 2 \times 8$ para $24 \times 3 \times 12$ ($h$ dividido por 1,5 nas três
    direções), o de $\lambda_1$ cai 5,0 vezes com as duas regras ($1{,}5^4 = 5{,}06$), uma inclinação de 1,43 a 1,44
    em $N$ (as equações crescem 3,06 vezes, menos que $1{,}5^3$, pelos nós do contorno).
  \item A malha $24 \times 3 \times 12$ com $5^3$ pontos (8575 equações) fica a $3{,}3 \times 10^{-5}$ da referência
    nos 20 modos ($\lambda_1$: $1{,}2 \times 10^{-6}$), no nível de Galerkin 2 (6642 graus de liberdade:
    $2{,}7 \times 10^{-5}$ e $9{,}7 \times 10^{-7}$); Galerkin 1 (1050) fica a $4{,}6 \times 10^{-4}$ e
    $1{,}7 \times 10^{-5}$. Galerkin 1 e 2 também ficam acima da referência nos 20 modos, de
    $9{,}7 \times 10^{-7}$ ($\lambda_1$) a $2{,}7 \times 10^{-5}$ ($\lambda_{20}$) em Galerkin 2.
  \item A soma de todos os autovalores (o traço de $\mathbf{B}^{-1}\mathbf{C}$, a integral em $D$ da diagonal do
    núcleo projetado no espaço da malha) fica abaixo de $270\pi$ e sobe com a malha: 0,9489 a 0,9782 de $270\pi$ com
    $3^3$ pontos e 0,8956 a 0,9634 com $5^3$. Na mesma malha, a regra de $3^3$ pontos dá o traço maior.
  \item Custo: o problema de autovalores denso domina (337 a 339~s e 1,7~GB com 8575 equações, todos os autovalores);
    a montagem de $\mathbf{C}$ leva 6,5~s ($3^3$) ou 78~s ($5^3$) com 2 threads. Galerkin 2 levou 17,3~h de formação
    e montagem num núcleo (máquinas diferentes; os tempos não são diretamente comparáveis).
\end{itemize}

\begin{figure}[htbp]
  \centering
  \includegraphics[width=\linewidth]{kl_cil_convergencia.pdf}
  \caption{Erro relativo dos autovalores do NeoPZ (Q2, geometria exata) contra a referência de Nyström, em função do
    número de equações: máximo nos 20 maiores modos (linhas cheias) e $\lambda_1$ (tracejadas), com $3^3$ e $5^3$
    pontos de Gauss por elemento em $\mathbf{C}$. Linhas pontilhadas: o erro máximo de Galerkin 1 e Galerkin 2 de
    Mika et al.\ (2021, Tabela~4); as estrelas marcam os seus graus de liberdade. A reta tem a inclinação de $h^4$
    ($N \propto h^{-3}$).}
  \label{fig:kl:cil:convergencia}
\end{figure}

\begin{table}[htbp]
  \centering
  \small
  \caption{Meio cilindro: erro relativo contra a referência de Nyström (máximo nos 20 maiores modos e em
    $\lambda_1$; todos positivos), soma de todos os autovalores, tempos de montagem de $\mathbf{C}$ (2 threads) e do
    problema de autovalores (todos, 1 thread) e pico de memória. Q2, geometria exata.}
  \label{tab:kl:cil:convergencia}
  \setlength{\tabcolsep}{3.5pt}
  \begin{tabular}{c r r c c c c r r r}
\toprule
malha & elem. & $N$ & pontos & máx. erro & erro & $\sum\lambda/(270\pi)$ & $\mathbf{C}$ (s) & aut. (s) & mem. (MB) \\
$n_\theta\times n_\rho\times n_z$ & & & & $k \le 20$ & $\lambda_1$ & & & & \\
\midrule
8$\times$1$\times$4 & 32 & 459 & $3^3$ & $1{,}8\cdot10^{-2}$ & $5{,}8\cdot10^{-4}$ & 0,9489 & 0,01 & 0,06 & 30 \\
12$\times$2$\times$6 & 144 & 1625 & $3^3$ & $3{,}0\cdot10^{-3}$ & $1{,}0\cdot10^{-4}$ & 0,9607 & 0,16 & 1,57 & 88 \\
16$\times$2$\times$8 & 256 & 2805 & $3^3$ & $1{,}1\cdot10^{-3}$ & $3{,}9\cdot10^{-5}$ & 0,9690 & 0,56 & 7,93 & 227 \\
24$\times$3$\times$12 & 864 & 8575 & $3^3$ & $2{,}2\cdot10^{-4}$ & $7{,}8\cdot10^{-6}$ & 0,9782 & 6,49 & 336,98 & 1724 \\
\midrule
8$\times$1$\times$4 & 32 & 459 & $5^3$ & $2{,}0\cdot10^{-3}$ & $9{,}6\cdot10^{-5}$ & 0,8956 & 0,11 & 0,04 & 30 \\
12$\times$2$\times$6 & 144 & 1625 & $5^3$ & $3{,}8\cdot10^{-4}$ & $1{,}7\cdot10^{-5}$ & 0,9301 & 2,12 & 1,54 & 91 \\
16$\times$2$\times$8 & 256 & 2805 & $5^3$ & $1{,}6\cdot10^{-4}$ & $6{,}2\cdot10^{-6}$ & 0,9458 & 6,67 & 7,92 & 231 \\
24$\times$3$\times$12 & 864 & 8575 & $5^3$ & $3{,}3\cdot10^{-5}$ & $1{,}2\cdot10^{-6}$ & 0,9634 & 78,24 & 338,92 & 1745 \\
\bottomrule
\end{tabular}

\end{table}

\begin{table}[htbp]
  \centering
  \small
  \caption{Os 20 maiores autovalores no meio cilindro: Galerkin 1 e Galerkin 2 de Mika et al.\ (2021, Tabela~4),
    NeoPZ na malha $24 \times 3 \times 12$ com $5^3$ pontos de Gauss e a referência de Nyström, e o erro relativo de
    cada um contra a referência.}
  \label{tab:kl:cil:autovalores}
  \begin{tabular}{r rrrr rrr}
\toprule
 & \multicolumn{4}{c}{$\lambda_k$} & \multicolumn{3}{c}{$(\lambda_k^h - \lambda_k)/\lambda_k$} \\
\cmidrule(lr){2-5}\cmidrule(lr){6-8}
$k$ & Galerkin 1 & Galerkin 2 & NeoPZ & Nyström & Galerkin 1 & Galerkin 2 & NeoPZ \\
\midrule
1 & 162,799154 & 162,796579 & 162,796625 & 162,796422 & $1{,}7\cdot10^{-5}$ & $9{,}7\cdot10^{-7}$ & $1{,}2\cdot10^{-6}$ \\
2 & 91,430631 & 91,428041 & 91,428082 & 91,427879 & $3{,}0\cdot10^{-5}$ & $1{,}8\cdot10^{-6}$ & $2{,}2\cdot10^{-6}$ \\
3 & 57,567029 & 57,564477 & 57,564521 & 57,564320 & $4{,}7\cdot10^{-5}$ & $2{,}7\cdot10^{-6}$ & $3{,}5\cdot10^{-6}$ \\
4 & 51,090179 & 51,087623 & 51,087676 & 51,087472 & $5{,}3\cdot10^{-5}$ & $3{,}0\cdot10^{-6}$ & $4{,}0\cdot10^{-6}$ \\
5 & 38,797409 & 38,794834 & 38,794874 & 38,794672 & $7{,}1\cdot10^{-5}$ & $4{,}2\cdot10^{-6}$ & $5{,}2\cdot10^{-6}$ \\
6 & 27,903861 & 27,901284 & 27,901329 & 27,901128 & $9{,}8\cdot10^{-5}$ & $5{,}6\cdot10^{-6}$ & $7{,}2\cdot10^{-6}$ \\
7 & 25,056111 & 25,053562 & 25,053612 & 25,053410 & $1{,}1\cdot10^{-4}$ & $6{,}0\cdot10^{-6}$ & $8{,}1\cdot10^{-6}$ \\
8 & 19,368934 & 19,366594 & 19,366627 & 19,366437 & $1{,}3\cdot10^{-4}$ & $8{,}1\cdot10^{-6}$ & $9{,}8\cdot10^{-6}$ \\
9 & 16,157004 & 16,154431 & 16,154470 & 16,154275 & $1{,}7\cdot10^{-4}$ & $9{,}6\cdot10^{-6}$ & $1{,}2\cdot10^{-5}$ \\
10 & 15,795308 & 15,792732 & 15,792776 & 15,792576 & $1{,}7\cdot10^{-4}$ & $9{,}9\cdot10^{-6}$ & $1{,}3\cdot10^{-5}$ \\
11 & 15,146229 & 15,143850 & 15,143879 & 15,143688 & $1{,}7\cdot10^{-4}$ & $1{,}1\cdot10^{-5}$ & $1{,}3\cdot10^{-5}$ \\
12 & 11,213287 & 11,210908 & 11,210949 & 11,210756 & $2{,}3\cdot10^{-4}$ & $1{,}4\cdot10^{-5}$ & $1{,}7\cdot10^{-5}$ \\
13 & 10,178969 & 10,176390 & 10,176431 & 10,176235 & $2{,}7\cdot10^{-4}$ & $1{,}5\cdot10^{-5}$ & $1{,}9\cdot10^{-5}$ \\
14 & 9,693564 & 9,690982 & 9,691009 & 9,690825 & $2{,}8\cdot10^{-4}$ & $1{,}6\cdot10^{-5}$ & $1{,}9\cdot10^{-5}$ \\
15 & 8,054783 & 8,052352 & 8,052390 & 8,052196 & $3{,}2\cdot10^{-4}$ & $1{,}9\cdot10^{-5}$ & $2{,}4\cdot10^{-5}$ \\
16 & 7,578086 & 7,576621 & 7,576627 & 7,576465 & $2{,}1\cdot10^{-4}$ & $2{,}1\cdot10^{-5}$ & $2{,}1\cdot10^{-5}$ \\
17 & 6,722378 & 6,719926 & 6,719953 & 6,719768 & $3{,}9\cdot10^{-4}$ & $2{,}4\cdot10^{-5}$ & $2{,}8\cdot10^{-5}$ \\
18 & 6,445577 & 6,443915 & 6,443922 & 6,443754 & $2{,}8\cdot10^{-4}$ & $2{,}5\cdot10^{-5}$ & $2{,}6\cdot10^{-5}$ \\
19 & 6,177345 & 6,174771 & 6,174786 & 6,174615 & $4{,}4\cdot10^{-4}$ & $2{,}5\cdot10^{-5}$ & $2{,}8\cdot10^{-5}$ \\
20 & 5,763787 & 5,761319 & 5,761355 & 5,761164 & $4{,}6\cdot10^{-4}$ & $2{,}7\cdot10^{-5}$ & $3{,}3\cdot10^{-5}$ \\
\midrule
g.l. & 1050 & 6642 & 8575 & -- & & & \\
\bottomrule
\end{tabular}

\end{table}

\begin{figure}[htbp]
  \centering
  \includegraphics[width=\linewidth]{kl_cil_espectro.pdf}
  \caption{Esquerda: os 20 maiores autovalores (Nyström; Galerkin 1 e 2 de Mika et al.; NeoPZ $24 \times 3 \times 12$
    com $5^3$ pontos; valores guardados no notebook \texttt{randomfied-v2.nb}). Direita: diferença relativa de cada um
    para a referência de Nyström do cilindro. O notebook usa o semi-anel poligonal, cujos autovalores exatos ficam
    0,35\% a 1,2\% abaixo dos do cilindro (Seção~\ref{sec:kl:cil:notebook}).}
  \label{fig:kl:cil:espectro}
\end{figure}

\subsection{O notebook \texttt{randomfied-v2.nb}}\label{sec:kl:cil:notebook}

\begin{sloppypar}
O notebook (Mathematica) monta $\mathbf{B}$ e $\mathbf{C}$ com hexaedros retos de 20 nós (\emph{serendipity} do
pacote \texttt{NDSolve`FEM}, com \texttt{MeshOrderAlteration} de ordem 2) numa malha de $16 \times 2 \times 4$ nós em
$\theta \times \rho \times z$ ($15 \times 1 \times 3$ elementos), com os nós em $(\rho\cos\theta, \rho\sin\theta, z)$:
o domínio é o semi-anel poligonal, de volume 842,0423 (contra $270\pi = 848{,}2300$). $\mathbf{C}$ usa
\texttt{ElementIntegrationPoints[HexahedronElement, 5]} em cada elemento do par e $\mathbf{B}$ a de 6. A mesma malha foi
calculada no NeoPZ (\texttt{GEOMETRIA=reta MALHAS=15x1x3}: 45 hexaedros Q2, 651 equações) e, de forma independente, em
numpy/scipy (\texttt{confere\_notebook.py}), em Q2 e no \emph{serendipity} de 20 nós do notebook (408 equações). A
referência é o Nyström refeito no semi-anel poligonal (os 15 hexaedros retos, só a simetria $z \to 15 - z$; as
estimativas extrapoladas diferem em até $6{,}2 \times 10^{-6}$ do valor). A Tabela~\ref{tab:kl:cil:notebook} resume:
\end{sloppypar}
\begin{itemize}
  \item o Q2 do numpy e o do NeoPZ dão os mesmos 20 autovalores com $3^3$ e $5^3$ pontos (diferença máxima de
    $4{,}9 \times 10^{-8}$, o arredondamento da 7.\textsuperscript{a} casa) e a mesma soma (804,0538099 e
    755,2080374); o \emph{serendipity} difere do Q2 em até $4{,}5 \times 10^{-4}$ (relativo);
  \item contra o Nyström do semi-anel, o NeoPZ com $5^3$ pontos fica de $-1{,}1 \times 10^{-2}$ a
    $+2{,}8 \times 10^{-3}$ (abaixo só nos modos 16 e 18); os valores guardados no notebook ficam de
    $-2{,}0 \times 10^{-2}$ a $+5{,}5 \times 10^{-1}$ (modo 9: 24,83 contra 15,97), e os seus modos 8 e 9 (24,8308 e
    24,8281) são quase iguais;
  \item o notebook guarda, logo depois dos autovalores, \texttt{Length[vec[[1]]]} = 92: os autovetores guardados têm
    92 componentes, enquanto a malha $15 \times 1 \times 3$ do código guardado tem 408 nós. A malha de
    $8 \times 2 \times 2$ nós ($7 \times 1 \times 1$ elementos de 20 nós) tem 92 nós, mas também não reproduz os
    valores ($\lambda_1 = 160{,}33$ com $5^3$ pontos em $\mathbf{C}$ e 167,48 com $2^3$).
\end{itemize}
Os autovalores guardados no notebook não saem do código guardado nele e não devem ser usados como referência; refeita,
a conta do notebook concorda com o NeoPZ.

\begin{table}[htbp]
  \centering
  \small
  \setlength{\tabcolsep}{4pt}
  \caption{Malha do notebook ($15 \times 1 \times 3$ hexaedros retos, semi-anel poligonal): referência de Nyström no
    semi-anel, Q2 do NeoPZ ($3^3$ e $5^3$ pontos em $\mathbf{C}$), Q2 do numpy ($5^3$), \emph{serendipity} de 20 nós
    (numpy, $5^3$), os valores guardados no notebook e os erros relativos contra a referência. Os autovetores guardados
    no notebook têm 92 componentes.}
  \label{tab:kl:cil:notebook}
  \begin{tabular}{r r rrr r r rr}
\toprule
 & Nyström & \multicolumn{3}{c}{Q2 (27 nós)} & serendipity & notebook & \multicolumn{2}{c}{$(\lambda_k^h - \lambda_k)/\lambda_k$} \\
\cmidrule(lr){3-5}\cmidrule(lr){8-9}
$k$ & (semi-anel) & NeoPZ $3^3$ & NeoPZ $5^3$ & numpy $5^3$ & 20 nós, $5^3$ & (guardado) & NeoPZ $5^3$ & notebook \\
\midrule
1 & 162,0687 & 162,1780 & 162,0871 & 162,0871 & 162,0871 & 163,9705 & $1{,}1\cdot10^{-4}$ & $1{,}2\cdot10^{-2}$ \\
2 & 90,7761 & 90,8849 & 90,7945 & 90,7945 & 90,7945 & 92,8766 & $2{,}0\cdot10^{-4}$ & $2{,}3\cdot10^{-2}$ \\
3 & 57,3273 & 57,4413 & 57,3425 & 57,3425 & 57,3424 & 56,1743 & $2{,}7\cdot10^{-4}$ & $-2{,}0\cdot10^{-2}$ \\
4 & 50,6386 & 50,7464 & 50,6569 & 50,6569 & 50,6569 & 53,2844 & $3{,}6\cdot10^{-4}$ & $5{,}2\cdot10^{-2}$ \\
5 & 38,5583 & 38,6733 & 38,5742 & 38,5742 & 38,5740 & 37,8918 & $4{,}1\cdot10^{-4}$ & $-1{,}7\cdot10^{-2}$ \\
6 & 27,6118 & 27,7185 & 27,6300 & 27,6300 & 27,6300 & 31,2591 & $6{,}6\cdot10^{-4}$ & $1{,}3\cdot10^{-1}$ \\
7 & 24,8632 & 24,9783 & 24,8797 & 24,8797 & 24,8795 & 30,6434 & $6{,}6\cdot10^{-4}$ & $2{,}3\cdot10^{-1}$ \\
8 & 19,2931 & 19,4159 & 19,2944 & 19,2944 & 19,2936 & 24,8308 & $7{,}2\cdot10^{-5}$ & $2{,}9\cdot10^{-1}$ \\
9 & 15,9718 & 16,0774 & 15,9898 & 15,9898 & 15,9898 & 24,8281 & $1{,}1\cdot10^{-3}$ & $5{,}5\cdot10^{-1}$ \\
10 & 15,6507 & 15,7655 & 15,6677 & 15,6677 & 15,6675 & 19,4531 & $1{,}1\cdot10^{-3}$ & $2{,}4\cdot10^{-1}$ \\
11 & 15,0666 & 15,1969 & 15,0711 & 15,0711 & 15,0700 & 18,2605 & $3{,}0\cdot10^{-4}$ & $2{,}1\cdot10^{-1}$ \\
12 & 11,1397 & 11,2759 & 11,1479 & 11,1479 & 11,1464 & 16,1117 & $7{,}3\cdot10^{-4}$ & $4{,}5\cdot10^{-1}$ \\
13 & 10,0747 & 10,1882 & 10,0918 & 10,0918 & 10,0915 & 13,6925 & $1{,}7\cdot10^{-3}$ & $3{,}6\cdot10^{-1}$ \\
14 & 9,5735 & 9,6777 & 9,5913 & 9,5913 & 9,5912 & 12,2071 & $1{,}9\cdot10^{-3}$ & $2{,}8\cdot10^{-1}$ \\
15 & 7,9915 & 8,1319 & 8,0031 & 8,0031 & 8,0011 & 10,8495 & $1{,}5\cdot10^{-3}$ & $3{,}6\cdot10^{-1}$ \\
16 & 7,5500 & 7,4890 & 7,4691 & 7,4691 & 7,4690 & 9,7552 & $-1{,}1\cdot10^{-2}$ & $2{,}9\cdot10^{-1}$ \\
17 & 6,6541 & 6,7474 & 6,6586 & 6,6586 & 6,6583 & 8,3436 & $6{,}9\cdot10^{-4}$ & $2{,}5\cdot10^{-1}$ \\
18 & 6,4088 & 6,3795 & 6,3464 & 6,3464 & 6,3463 & 7,5551 & $-9{,}7\cdot10^{-3}$ & $1{,}8\cdot10^{-1}$ \\
19 & 6,0982 & 6,2027 & 6,1155 & 6,1155 & 6,1153 & 6,9425 & $2{,}8\cdot10^{-3}$ & $1{,}4\cdot10^{-1}$ \\
20 & 5,7109 & 5,8516 & 5,7254 & 5,7254 & 5,7229 & 6,3018 & $2{,}6\cdot10^{-3}$ & $1{,}0\cdot10^{-1}$ \\
\midrule
$N$ & -- & 651 & 651 & 651 & 408 & -- & & \\
\bottomrule
\end{tabular}

\end{table}
